% !TEX root = ../main_en.tex

\section{Summary and Discussion}\label{sec:conclusion}

\subsection{Main Conclusions}

This document systematically compiles the theoretical analysis by Boltz and Ihle on density fluctuation suppression in antialigning active matter systems. The core findings include:

\begin{enumerate}
    \item \textbf{One-dimensional analytical solution}: Through the Poisson representation, the structure factor for a one-dimensional active lattice gas model with antialigning interactions is rigorously derived, revealing that long-range density fluctuations are significantly suppressed to approximately 1/2 of the Poisson value
    
    \item \textbf{Physical meaning of imaginary noise}: The imaginary noise that naturally emerges in the Poisson representation marks non-Poissonian statistics---this is the mathematical manifestation of the deviation of active systems from equilibrium statistics
    
    \item \textbf{Two-dimensional numerical verification}: Apparent hyperuniform behavior is observed in a Vicsek-like model, but it is emphasized that this may be a finite-size effect or a misinterpretation of a finite offset
    
    \item \textbf{Warning about numerical inference}: Without a theoretical basis, inferring hyperuniformity from numerical data is unreliable
\end{enumerate}

\subsection{Value of the Theoretical Method}

Advantages of the Poisson representation in this study:
\begin{itemize}
    \item Direct, non-perturbative access to fluctuations (the noise term in the final Langevin equation)
    \item Direct access to the correlations of the noise, which is crucial for determining the structure factor
    \item In contrast to certain path integral methods (which asymptotically derive Gaussian uncorrelated noise)
\end{itemize}

\subsection{Open Questions}

\begin{enumerate}
    \item \textbf{Strict hyperuniformity}: Does true hyperuniformity ($S(0) = 0$) genuinely exist in two-dimensional continuous systems, or is it merely $k^2$ behavior with a finite offset?
    \item \textbf{Higher-dimensional generalization}: How does the analytical description of the one-dimensional lattice model generalize to higher dimensions?
    \item \textbf{Experimental connections}: Comparison with experiments on antialigning self-driven particles (such as the works cited in\cite{boltz2025})
    \item \textbf{Tuning of point pattern statistics}: Antialigning self-driven particles as an independent dynamical process for generating disordered point patterns with tunable statistics
\end{enumerate}

\begin{notebox}[Final Remarks]
The transition from Poisson statistics at $v = 0$ (passive system) to non-Poissonian, possibly hyperuniform statistics at $v > 0$ (active system) is a remarkable manifestation of the relevance of long-range correlations even in the absence of direct order.

This suppression of density fluctuations is purely a result of dynamically established correlations (due to aligning interactions), rather than more traditional mechanisms such as elasticity. It is precisely this long-range effect that enables the one-dimensional analytically solvable toy model to capture and connect with the more complex two-dimensional Vicsek-like model.
\end{notebox}
